% فصل سوم: شبکه‌های عصبی آگاه از فیزیک (روش پژوهش)
\chapter{شبکه‌های عصبی آگاه از فیزیک}
\label{ch:method}

در این فصل، ساختار ریاضی و الگوریتمی چارچوب «شبکه‌های عصبی آگاه از فیزیک» (\lr{PINN}) که متدولوژی اصلی این پژوهش را تشکیل می‌دهد، به‌تفصیل و به‌صورتی گام‌به‌گام و خودخوان تشریح می‌گردد. هدف این چارچوب، ایجاد پلی استوار میان قوانین پایه‌ای فیزیک و توانمندی تقریب‌زدن شبکه‌های عصبی عمیق است. در ادامه، ابتدا صورت‌بندی کلی مسئله و ایده محوری روش تبیین شده و سپس ساختار مدل‌های زمان‌پیوسته و زمان‌گسسته در مسائل مستقیم معرفی می‌شوند. پس از آن، فرمول‌بندی مسئله معکوس برای کشف پارامترهای مجهول معادلات ارائه می‌گردد و در پایان، ملاحظات تجربی و عملیاتی حاکم بر آموزش این مدل‌ها مورد بحث قرار می‌گیرد \cite{raissi2019}.

\section{صورت‌بندی کلی مسئله و ایده محوری}
\label{sec:formulation}

بسیاری از سیستم‌های دینامیکی در فیزیک و مهندسی توسط معادلات دیفرانسیل جزئی غیرخطی و وابسته به زمان توصیف می‌شوند. صورت پارامتری عام این معادلات را می‌توان به فرم زیر بیان کرد:
\begin{equation}
u_t + \mathcal{N}[u; \lambda] = 0, \quad \mathbf{x} \in \Omega, \quad t \in [0,T],
\label{eq:general-pde}
\end{equation}
که در آن $u(t,\mathbf{x})$ تابع مجهول حالت سیستم (مانند سرعت، فشار، دما یا غلظت)، $u_t = \frac{\partial u}{\partial t}$ مشتق زمانی متغیر حالت، $\mathcal{N}[\cdot; \lambda]$ عملگر دیفرانسیلی غیرخطی شامل مشتقات مکانی و جملات غیرخطی وابسته به بردار پارامترهای فیزیکی $\lambda \in \mathbb{R}^p$، و $\Omega \subset \mathbb{R}^d$ دامنه مکانی پیوسته با مرز $\partial\Omega$ است. این رابطه بیانگر تعادل نیروها و قوانین بقای حاکم بر فیزیک سامانه در هر نقطه مکانی و در هر لحظه زمانی است. به‌عنوان نمونه، در معادله برگرز با لزجت، عملگر دیفرانسیلی به‌صورت $\mathcal{N}[u;\lambda] = \lambda_1 u u_x - \lambda_2 u_{xx}$ با بردار پارامتر $\lambda = (\lambda_1, \lambda_2)$ تعریف می‌گردد.

با فرض در اختیار داشتن مجموعه‌ای از داده‌های اندازه‌گیری‌شده (که ممکن است محدود، نامنظم یا آلوده به نویز باشند)، دو مسئله بنیادین قابل طرح است:
\begin{enumerate}
\item \textbf{مسئله مستقیم (\lr{Forward Problem}):} با فرض معلوم بودن کامل بردار پارامترها $\lambda$ و شرایط مرزی و اولیه، هدف محاسبه و بازسازی میدان پاسخ $u(t,\mathbf{x})$ در سراسر دامنه زمانی-مکانی است. در این مسئله، شبکه عصبی نقش یک حل‌گر عددی بدون شبکه را ایفا می‌کند.
\item \textbf{مسئله معکوس (\lr{Inverse Problem}):} در شرایطی که بردار پارامترهای $\lambda$ مجهول بوده و مقادیری از متغیر حالت $u(t,\mathbf{x})$ در نقاط پراکنده‌ای از دامنه اندازه‌گیری شده باشد، هدف کشف و تخمین دقیق پارامترهای فیزیکی $\lambda$ هم‌زمان با برازش میدان پاسخ است. در این حالت، شبکه عصبی فرآیند شناسایی سیستم را به انجام می‌رساند.
\end{enumerate}

ایده محوری در شبکه‌های آگاه از فیزیک بر پایه تقریب تابع جواب $u(t,\mathbf{x})$ توسط یک شبکه عصبی عمیق $u_\theta(t,\mathbf{x})$ با بردار پارامترهای وزن و بایاس $\theta$ بنا شده است. بر این اساس، عملگر باقیمانده دیفرانسیلی فیزیک $f(t,\mathbf{x})$ به‌صورت زیر تعریف می‌گردد:
\begin{equation}
f(t,\mathbf{x}) := u_t + \mathcal{N}[u; \lambda],
\label{eq:residual}
\end{equation}
که در آن $u = u_\theta(t,\mathbf{x})$ خروجی شبکه عصبی است. ویژگی برجسته این چارچوب آن است که کلیه مشتقات جزئی ظاهرشده در $f$ (نظیر $u_t$ و مشتقات مکانی) با استفاده از مشتق‌گیری خودکار نسبت به ورودی‌های گراف محاسباتی ارزیابی می‌شوند. در نتیجه، $f(t,\mathbf{x})$ خود نمایانگر یک شبکه عصبی با پارامترهای مشترک $\theta$ است که اصطلاحاً «شبکه باقیمانده آگاه از فیزیک» نامیده می‌شود.

در صورتی که شبکه $u_\theta$ به جواب دقیق معادله میل کند، مقدار باقیمانده در تمامی نقاط دامنه به سمت صفر میل خواهد کرد ($f \to 0$). بنابراین، فرایند آموزش از طریق یک تابع هزینه ترکیبی که هم‌زمان خطای برازش داده‌ها و خطای عدم ارضای معادله دیفرانسیل را کمینه می‌سازد، هدایت می‌شود. شکل \ref{fig:pinn} نمای شماتیک این ساختار یکپارچه را نشان می‌دهد.

\begin{figure}[t]
\centering
\begin{latin}
\begin{tikzpicture}[>=Stealth,
  block/.style={rectangle, draw=myblue!80!black, fill=myblue!5, thick, rounded corners=2pt, minimum width=2.5cm, minimum height=0.95cm, align=center, font=\footnotesize\bfseries},
  procbox/.style={rectangle, draw=black!75, fill=black!4, thick, rounded corners=2pt, minimum width=2.2cm, minimum height=0.95cm, align=center, font=\footnotesize},
  lossbox/.style={rectangle, draw=red!70!black, fill=red!5, thick, rounded corners=2pt, minimum width=2.8cm, minimum height=0.95cm, align=center, font=\footnotesize\bfseries},
  optbox/.style={rectangle, draw=myblue!80!black, fill=myblue!10, thick, rounded corners=2pt, minimum width=2.4cm, minimum height=0.95cm, align=center, font=\footnotesize\bfseries}
]
  % Top row: Input -> Neural Network -> Autodiff -> PDE Residual
  \node[procbox] (input) at (0, 0) {Input\\$(t, \mathbf{x})$};
  \node[block] (net) at (3.2, 0) {Neural Network\\$u_\theta(t, \mathbf{x})$};
  \node[procbox] (ad) at (6.6, 0) {Automatic\\Diff ($\partial_t, \partial_x, \dots$)};
  \node[block] (res) at (10.0, 0) {PDE Residual\\$f_\theta(t, \mathbf{x})$};

  % Bottom row: Optimizer and Composite Loss
  \node[optbox] (opt) at (3.2, -2.3) {Optimizer\\(Adam / L-BFGS)};
  \node[lossbox] (loss) at (8.3, -2.3) {Composite Loss\\$\mathrm{MSE}_u + \mathrm{MSE}_f$};

  % Top row arrows
  \draw[->, thick] (input) -- (net);
  \draw[->, thick] (net) -- (ad);
  \draw[->, thick] (ad) -- (res);

  % Downward arrows to Loss
  \draw[->, thick] (net.south) -- ++(0,-0.5) -| ([xshift=-0.6cm]loss.north) 
    node[pos=0.25, left, font=\scriptsize] {$u$ on IC/BC};
  \draw[->, thick] (res.south) -- ++(0,-0.5) -| ([xshift=0.6cm]loss.north) 
    node[pos=0.25, right, font=\scriptsize] {$f$ on Collocation};

  % Loss to Optimizer
  \draw[->, thick] (loss) -- (opt) node[midway, above, font=\scriptsize] {$\nabla_\theta \mathcal{L}$};

  % Feedback: Optimizer updates Network parameters theta
  \draw[->, thick, dashed, myblue!80!black] (opt.north) -- (net.south) 
    node[midway, right, font=\scriptsize] {Update $\theta$};

\end{tikzpicture}
\end{latin}
\caption{نمودار بلوکی چارچوب شبکه عصبی آگاه از فیزیک: تعامل شبکه پاسخ، مشتق‌گیری خودکار، شبکه باقیمانده و حلقه بهینه‌سازی}
\label{fig:pinn}
\end{figure}

\section{مدل‌های زمان‌پیوسته برای مسئله مستقیم}
\label{sec:continuous}

در مدل‌های زمان‌پیوسته، متغیرهای فضا و زمان به‌صورت هم‌زمان به‌عنوان ورودی‌های پیوسته به شبکه عصبی تغذیه می‌شوند: $(t,\mathbf{x}) \mapsto u_\theta(t,\mathbf{x})$. این ویژگی موجب می‌شود که شبکه به یک تقریب‌زننده پیوسته در تمام دامنه فضا-زمان بدل گردد. پارامترهای مجهول شبکه $\theta$ با کمینه‌سازی تابع هزینه ترکیبی زیر بهینه‌سازی می‌شوند:
\begin{equation}
\mathcal{L}(\theta) = \mathrm{MSE}_u(\theta) + \mathrm{MSE}_f(\theta),
\label{eq:loss-forward}
\end{equation}
که در آن:
\begin{equation}
\begin{aligned}
\mathrm{MSE}_u(\theta) &= \frac{1}{N_u}\sum_{i=1}^{N_u} \left|u_\theta(t^i_u,\mathbf{x}^i_u) - u^i\right|^2, \\
\mathrm{MSE}_f(\theta) &= \frac{1}{N_f}\sum_{i=1}^{N_f} \left|f_\theta(t^i_f,\mathbf{x}^i_f)\right|^2.
\end{aligned}
\label{eq:loss-terms}
\end{equation}
در روابط فوق، مجموعه $\{t^i_u, \mathbf{x}^i_u, u^i\}_{i=1}^{N_u}$ شامل نقاط داده روی شرایط مرزی و اولیه است که مقادیر دقیق متغیر حالت در آن‌ها مشخص است. مجموعه $\{t^i_f, \mathbf{x}^i_f\}_{i=1}^{N_f}$ نمایانگر «نقاط هم‌مکانی»\LTRfootnote{Collocation Points} است؛ نقاطی در درون دامنه پیوسته زمانی-مکانی که هیچ اندازه‌گیری تجربی از مقدار $u$ در آن‌ها وجود ندارد و صرفاً ارضای ساختار معادله دیفرانسیل ($f \approx 0$) بر آن‌ها اعمال می‌گردد.

جمله $\mathrm{MSE}_f$ در حقیقت نقش یک منظم‌ساز فیزیکی قدرتمند را ایفا می‌کند که فضای فرضیات شبکه عصبی را به توابعی محدود می‌سازد که ساختار فیزیکی حاکم را برآورده کنند. در مسئله مستقیم معادله برگرز با لزجت $\nu$، باقیمانده فیزیکی صریح به صورت زیر تشکیل می‌گردد:
\begin{equation}
f := u_t + u u_x - \nu u_{xx}.
\label{eq:burgers-f}
\end{equation}
این ساختار به مدل اجازه می‌دهد با تعداد بسیار اندکی از داده‌های مرزی-اولیه، میدان پیوسته پاسخ را در سراسر دامنه بازسازی نماید.

\section{مدل‌های زمان‌گسسته}
\label{sec:discrete}

در مسائلی که نیازمند پیشروی در بازه‌های زمانی بسیار طولانی یا دقت‌های زمانی فوق‌العاده بالا هستند، مدل‌های زمان‌پیوسته ممکن است با افزایش تعداد نقاط هم‌مکانی و دشواری بهینه‌سازی روبه‌رو شوند. برای رفع این چالش، فرمول‌بندی زمان‌گسسته با تلفیق طرح‌های گام‌زنی زمانی ضمنی از خانواده روش‌های رانگه-کوتا\LTRfootnote{Runge-Kutta Methods} با $q$ مرحله میانی پیشنهاد شده است.

با اعمال صورت کلی یک روش رانگه-کوتای ضمنی به معادله دیفرانسیل \eqref{eq:general-pde}، روابط گام زمانی به صورت زیر نوشته می‌شوند:
\begin{equation}
\begin{aligned}
u^{n+c_i} &= u^n - \Delta t \sum_{j=1}^{q} a_{ij}\, \mathcal{N}[u^{n+c_j}], \quad i=1,\dots,q, \\
u^{n+1} &= u^n - \Delta t \sum_{j=1}^{q} b_j\, \mathcal{N}[u^{n+c_j}],
\end{aligned}
\label{eq:rk}
\end{equation}
که در آن $\{a_{ij}, b_j, c_j\}$ ضرایب جدول بوچر\LTRfootnote{Butcher Tableau} طرح رانگه-کوتا هستند \cite{butcher2016}. در این مدل، خروجی شبکه عصبی شامل $q+1$ مقدار است که متناظر با جواب در مراحل میانی $u^{n+c_1},\dots,u^{n+c_q}$ و جواب گام بعدی $u^{n+1}$ می‌باشد. تابع هزینه بر اساس تطابق این مقادیر با روابط \eqref{eq:rk} و داده‌های لحظات $t^n$ و $t^{n+1}$ شکل می‌گیرد.

مزیت عمده این رهیافت، پایداری ذاتی روش‌های ضمنی رانگه-کوتا (نظیر خانواده گاوس-لژاندر) است که امکان انتخاب گام‌های زمانی بزرگ ($\Delta t$) بدون نقض شرایط پایداری را فراهم می‌آورد. در این پژوهش، تمرکز پیاده‌سازی و ارزیابی تجربی بر مدل زمان‌پیوسته قرار دارد، اما تحلیل مبانی مدل‌های زمان‌گسسته برای درک جامع ابعاد چارچوب ضروری است.

\section{مسئله معکوس و کشف پارامترهای فیزیکی}
\label{sec:inverse}

در مسائل معکوس یا کشف معادلات حاکم، علاوه بر میدان مجهول $u(t,\mathbf{x})$، بردار ضرایب فیزیکی $\lambda$ نیز مجهول بوده و باید مستقیماً از روی مشاهدات تجربی استخراج گردد. در چارچوب شبکه‌های آگاه از فیزیک، این فرایند با تعریف $\lambda$ به‌عنوان متغیرهای بهینه‌سازی اضافی در کنار وزن‌های شبکه صورت می‌پذیرد.

برای نمونه، در مسئله کشف پارامترهای معادله برگرز، عملگر باقیمانده به شکل پارامتری زیر در نظر گرفته می‌شود:
\begin{equation}
f := u_t + \lambda_1 u u_x - \lambda_2 u_{xx},
\label{eq:burgers-f-inv}
\end{equation}
که در آن $\lambda_1$ و $\lambda_2$ پارامترهای اسکالر ناشناخته هستند. فرایند بهینه‌سازی تابع هزینه \eqref{eq:loss-forward} هم‌زمان نسبت به $(\theta, \lambda_1, \lambda_2)$ انجام می‌پذیرد.

تفاوت بنیادین این رهیافت با روش‌های متداول تحلیل داده (نظیر رگرسیون نمادین یا مشتق‌گیری تفاضلی روی داده‌های نویزی) در آن است که مشتقات جزئی مورد نیاز از روی شبکه هموار $u_\theta$ با مشتق‌گیری خودکار ارزیابی می‌شوند. این امر باعث می‌شود که شبکه عصبی هم‌زمان نقش یک فیلتر هموارکننده نویز را ایفا کرده و پایداری محاسباتی بالایی در برابر خطاهای اندازه‌گیری فراهم آورد.

\section{ملاحظات عملی در پیاده‌سازی و آموزش}
\label{sec:practical}

تجارب محاسباتی و مبانی تجربی حاکی از آن است که دستیابی به همگرایی مطلوب در شبکه‌های آگاه از فیزیک مستلزم رعایت چند اصل کلیدی در طراحی ساختار و الگوریتم آموزش است:

\textbf{معماری شبکه و توابع فعال‌ساز.} در اغلب مسائل فیزیکی، پرسپترون‌های چندلایه با عمق ۴ تا ۸ لایه و ۲۰ تا ۵۰ نورون در هر لایه با تابع فعال‌ساز تانژانت هیپربولیک ($\tanh$) عملکرد مطلوبی نشان می‌دهند. به‌کارگیری لایه‌های افت تصادفی (\lr{Dropout}) در این شبکه‌ها معمولاً غیرضروری است؛ زیرا قید فیزیکی باقیمانده به‌طور خودکار مانع بیش‌برازش می‌گردد.

\textbf{راهبرد بهینه‌سازی ترکیبی.} استفاده از الگوریتم دو مرحله‌ای شامل شروع با بهینه‌ساز آدام (جهت فرار از کمینه‌های محلی و کاهش سریع خطای اولیه) و سپس تغییر وضعیت به بهینه‌ساز شبه‌نیوتنی \lr{L-BFGS} (جهت رسیدن به دقت‌های بالا و همگرایی مرتبه دو) پایدارترین استراتژی آموزشی است.

\textbf{راهبرد نمونه‌برداری نقاط هم‌مکانی.} توزیع و تراکم نقاط هم‌مکانی در دامنه حل بر پایداری مدل اثرگذار است. در کنار نمونه‌برداری یکنواخت تصادفی یا روش ابرمکعب لاتین\LTRfootnote{Latin Hypercube Sampling}، افزایش چگالی نقاط هم‌مکانی در نواحی با تغییرات مکانی شدید (نظیر جبهه شوک در معادله برگرز) دقت تقریب را به میزان چشمگیری بهبود می‌بخشد.

جدول \ref{tab:models} مشخصات و تمایزات ساختاری مدل‌های زمان‌پیوسته و زمان‌گسسته را مقایسه می‌نماید.

\begin{table}[t]
\caption{مقایسه تطبیقی مدل‌های زمان‌پیوسته و زمان‌گسسته در چارچوب شبکه‌های آگاه از فیزیک}
\label{tab:models}
\centering
\singlespacing
\begin{tabular}{|p{2.8cm}|p{5.2cm}|p{5.6cm}|}
\hline
\textbf{شاخص ساختاری} & \textbf{مدل زمان‌پیوسته} & \textbf{مدل زمان‌گسسته} \\
\hline
ورودی شبکه & بردار پیوسته زمانی-مکانی $(t, \mathbf{x})$ & متغیر مکانی $\mathbf{x}$ در یک مقطع زمانی \\
\hline
خروجی شبکه & مقدار اسکالر میدان $u(t,\mathbf{x})$ & مقادیر مراحل میانی رانگه--کوتا و~$u^{n+1}$ \\
\hline
داده‌های آموزشی & شرایط اولیه و مرزی در دامنه زمان & مقادیر حالت در دو لحظه گسسته $t^n$ و $t^{n+1}$ \\
\hline
مزیت شاخص & سادگی و یکپارچگی ساختار & پایداری در گام‌های زمانی بزرگ \\
\hline
وضعیت در این پژوهش & پیاده‌سازی و تحلیل کامل تجربی & تحلیل و بررسی نظری \\
\hline
\end{tabular}
\end{table}

\section{جمع‌بندی فصل}
\label{sec:ch3-summary}

در این فصل، ساختار ریاضی و متدولوژی شبکه‌های عصبی آگاه از فیزیک تشریح شد: مکانیزم تعریف شبکه باقیمانده بر پایه مشتق‌گیری خودکار، طراحی توابع هزینه ترکیبی فیزیک-داده، فرمول‌بندی مدل‌های زمان‌پیوسته و زمان‌گسسته برای مسائل مستقیم و ساختار استخراج ضرایب در مسائل معکوس. در فصل چهارم، پیاده‌سازی این چارچوب برای معادله برگرز یک‌بعدی و نتایج عددی حاصل ارائه خواهد شد.
